\section{Results}
\label{sec:results}

We present results in the order of the paper's narrative: first, the
transcoder-guided method (\method{}) beats \baseline{} at matched budget
(\S\ref{sec:results-spider}, \S\ref{sec:results-downstream}); second, a
controlled ablation dissects \emph{why}, revealing the conjunction as the
active ingredient and the dictionary as optional (\S\ref{sec:ablation});
third, the bit-width analysis locates where any of this matters
(\S\ref{sec:results-3bit}).

\subsection{Spider text-to-SQL (primary head-to-head)}
\label{sec:results-spider}

Table~\ref{tab:spider} reports execution accuracy on the Spider dev set (1034
examples, test-suite-sql-eval) under the matched protocol. v2~mult averages
29.4\% over four paired calibration seeds vs.\ \baseline{} 26.5\%
($+2.9$\,pp mean paired $\Delta$); the best seed reaches 32.6\%. Replacing CE
entirely with transcoder alignment (v1) ties \baseline{} (25.0 vs.\ 26.0)---a
first hint, confirmed by the ablation below, that the improvement requires the
\emph{conjunction}, not the transcoder signal alone. At 3-bit, \baseline{}
reaches 64.8\%, v1 ties (65.0\%), and v2~mult reaches 66.2\% (seed 0)---again
a compressed gap at the milder bit-width.

\begin{table}[t]
\centering
\caption{Spider execution accuracy (\%) on Llama-3.1-8B-Instruct, 0.35\%
budget, 1360 Spider-train contrastive calibration pairs. Multi-seed: mean $\pm$
sample std over four calibration seeds.}
\label{tab:spider}
\begin{tabular}{lcc}
\toprule
Method & FP16 & 2-bit \\
\midrule
FP16 (unquantized)          & 67.8 & --   \\
\baseline{}~\citep{xiao2025tacq} (mean) & -- & 26.5 $\pm$ 2.0 \\
\method{} v1 (align-only)   & --   & 25.0 \\
\method{} v2 mult (mean)    & --   & \textbf{29.4 $\pm$ 2.7} \\
\method{} v2 mult (best seed) & -- & 32.6 \\
\bottomrule
\end{tabular}
\end{table}

\paragraph{Replication sanity check (Llama-3-8B).}
Our harness reproduces \baseline{} Spider exec accuracy: 67.7 (FP16), 22.3
(2-bit), 60.1 (3-bit) vs.\ paper 67.6 / 21.92 / 58.32~\citep{xiao2025tacq}.

\subsection{Task-conditioned downstream comparison (multi-seed)}
\label{sec:results-downstream}

Tables~\ref{tab:downstream2} and \ref{tab:downstream3} report the firm
head-to-head between \baseline{} and \method{} v2~mult under the
task-conditioned protocol (calibrate and evaluate on the same task;
\S\ref{sec:setup}), with seeds varying the contrastive calibration sample.

\begin{table}[t]
\centering
\caption{Task-conditioned \textbf{2-bit} (0.35\% budget): mean over calibration
seeds (seed list in parentheses). GSM8K: 8-shot CoT accuracy; MMLU: held-out
last 25\% of test. \todo{fill rebuilt multi-seed numbers when jobs land;
single-seed entries marked $^\dagger$.}}
\label{tab:downstream2}
\begin{tabular}{lccc}
\toprule
Task & FP16 & \baseline{} & \method{} v2 mult \\
\midrule
GSM8K                & 69.3 & 28.3$^\dagger$ & 29.8$^\dagger$ \\
MMLU STEM            & 59.0 & \todo{TBD} & \todo{TBD} \\
MMLU humanities      & 66.0 & \todo{TBD} & 47.5$^\dagger$ \\
MMLU social sciences & 75.6 & \todo{TBD} & \todo{TBD} \\
\bottomrule
\end{tabular}
\end{table}

\begin{table}[t]
\centering
\caption{Task-conditioned \textbf{3-bit}: same protocol. \baseline{} over seeds
$\{0,1,2\}$, v2 mult over $\{0,1\}$ (seed-2 in flight). All gaps are within
seed spread: at 3-bit the methods are statistically tied.}
\label{tab:downstream3}
\begin{tabular}{lccc}
\toprule
Task & FP16 & \baseline{} & \method{} v2 mult \\
\midrule
GSM8K                & 69.3 & 63.7 & 63.3 \\
MMLU STEM            & 59.0 & 56.5 & 54.8 \\
MMLU humanities      & 66.0 & 63.0 & 63.1 \\
MMLU social sciences & 75.6 & 73.0 & 74.1 \\
\bottomrule
\end{tabular}
\end{table}

\subsection{Saliency-source ablation: what carries the gain?}
\label{sec:ablation}

Having established that \method{} beats \baseline{} at matched budget, we ask
\emph{which ingredient does the work}. The ablation holds \emph{everything}
fixed---model (Llama-3.1-8B-Instruct), FP16 budget (0.35\%, global top-$k$),
contrastive calibration pairs (128), GPTQ (true-sequential), and
evaluation---and varies only the saliency signal that ranks weights. Eight
arms span the hypothesis space from task-agnostic controls to the full
conjunction, with each circuit arm instantiated in both the transcoder basis
and the model's native MLP-neuron basis (Table~\ref{tab:ablation}).

\begin{table}[t]
\centering
\caption{Saliency-source ablation at 2-bit (0.35\% FP16 budget, seed 0).
Identical pipeline in every row; only the weight-ranking signal differs.
GSM8K: 8-shot CoT accuracy on the full test set. MMLU-STEM: 5-shot accuracy on
the held-out last 25\% of each subject's test split (chance $=25\%$).
Dict.\ = requires an external feature dictionary.}
\label{tab:ablation}
\begin{tabular}{llccc}
\toprule
Arm & Saliency signal & Dict. & GSM8K & MMLU-STEM \\
\midrule
random    & uniform random                          & --  & 2.5  & 30.3 \\
weight    & $|W|$                                   & --  & 14.4 & 33.6 \\
magnitude & $|W| \cdot |\Delta W_{\text{quant}}|$   & --  & 16.7 & 27.8 \\
\midrule
ce ($\approx$\baseline{}) & CE gradient saliency    & --  & 25.4 & 41.3 \\
align        & circuit only, transcoder basis       & \checkmark & 25.6 & 39.2 \\
align$_{\text{free}}$ & circuit only, native basis  & --  & 22.9 & 42.6 \\
\midrule
mult (\method{})            & CE $\cap$ circuit, transcoder & \checkmark & \textbf{28.7} & 38.9 \\
mult$_{\text{free}}$ (\methodfree{}) & CE $\cap$ circuit, native & -- & 26.5 & \textbf{44.2} \\
\bottomrule
\end{tabular}
\end{table}

\paragraph{Finding 1: task-circuit information dominates bit-placement
heuristics.} Every circuit-informed arm beats every control by a wide margin.
On GSM8K, protecting the highest-magnitude weights recovers only 14.4\% and
magnitude$\times$error 16.7\%, versus 22.9--28.7\% for circuit-informed
selection; random placement collapses to 2.5\%. On MMLU-STEM the controls sit
near chance (27.8--33.6\%) while circuit arms reach 38.9--44.2\%. \emph{Where}
the task circuit lives carries most of the recoverable signal at 2-bit.

\paragraph{Finding 2: the conjunction is the active ingredient.}
On GSM8K the conjunction beats both of its parts: mult 28.7\% vs.\ ce 25.4\%
(loss-only) and align 25.6\% (circuit-only). On MMLU-STEM the same holds in the
native basis: mult$_{\text{free}}$ 44.2\% vs.\ ce 41.3\% and
align$_{\text{free}}$ 42.6\%. Neither signal alone suffices; the AND-gate over
task loss and task circuit is what buys accuracy. This retroactively explains
the Spider v1 result (\S\ref{sec:results-spider}): align-only tied \baseline{}
because it had discarded the loss half of the conjunction.

\paragraph{Finding 3: the dictionary is optional---and sometimes harmful.}
\methodfree{} recovers 92\% of the transcoder conjunction on GSM8K (26.5 vs.\
28.7) and \emph{surpasses every transcoder arm} on MMLU-STEM (44.2 vs.\ 38.9
for mult). We attribute the reversal to basis--task fit: the CRV transcoders
were trained to reconstruct general instruct-model behavior, and their feature
space may tile math-word-problem structure (GSM8K) better than the
heterogeneous knowledge retrieval MMLU requires; the native neuron basis is
task-neutral. Practically, the external dictionary is an optional booster on
some tasks rather than a dependency: conjunctive task-circuit PTQ works
out-of-the-box on any transformer.

\subsection{Conjunctive saliency is a low-bit phenomenon}
\label{sec:results-3bit}

Comparing Tables~\ref{tab:ablation} and \ref{tab:downstream3} yields the
paper's sharpest intuition. At \textbf{2-bit}, the saliency source decides
between collapse and function: the spread between the best and worst
task-conditioned arm is 3--6 points, and between conjunction and magnitude
controls 10--26 points. At \textbf{3-bit}, every gap between
\baseline{} and the conjunction shrinks below seed noise ($\leq$1.7 points,
Table~\ref{tab:downstream3}). The mechanism is simple: 3-bit quantization noise
is small enough that GPTQ's error compensation absorbs most task damage
regardless of which 0.35\% of weights are protected; at 2-bit the noise is
catastrophic and the FP16 budget must land on weights that are simultaneously
task-critical and quantization-fragile---precisely what the conjunction
selects. Task-circuit saliency therefore matters most exactly when the bit
budget is harshest, which is the regime mixed-precision PTQ exists for.

\paragraph{A cautionary note on bit-width bookkeeping.}
During this study we found that an earlier run had silently built ``2-bit''
\method{} engines at 3 bits (a quantizer default overriding the pipeline
intent), producing implausibly strong 2-bit numbers ($\sim$62\% GSM8K) that a
matched-budget ablation immediately exposed ($28.7\%$). We recommend
controlled saliency-source ablations as standard practice for mixed-precision
PTQ papers: they are cheap, and they catch both bookkeeping errors and
overclaiming.

\paragraph{Dispute sweep (anchored union).}
The anchored-union construction (\S\ref{sec:method-anchored}) locks the
consensus core ($\sim$46\% of budget at 2-bit) and re-ranks only the XOR
dispute mass. The mask morph behaves as designed ($J$ vs.\ \baseline{} falls
from 0.708 to 0.462 as $\lambda$ goes 0 to 1). Partial Spider evals at 3-bit
peak at the \baseline{}-heavy endpoint ($\lambda{=}0$: 63.2\% vs.\ 59.9\% at
$\lambda{=}0.5$), consistent with milder quantization favoring load-bearing
MLP mass over routing-heavy circuit weights. \todo{complete 2-bit sweep.}
